\section{\Nom{} World Models on Frozen DINO Tokens}
\label{sec:method}

\textbf{Setting.} A frozen encoder maps each observation to patch tokens $s_t \in \R^{N \times D}$ with $N = 256$ tokens on a $16\times16$ grid and $D = 384$ channels: DINOv2 ViT-S/14 at $224$\,px \citep{oquab2024dinov2} for the bounce, push, Push-T and BAIR benchmarks, and FlowWM's DINOv3 feature definition \citep{simeoni2025dinov3,porcher2026flowwm} for Bouncing Shapes. Tokens are precomputed and standardized per token. The model takes a context of three token maps $c_t$, an action $a_t$ and a noise vector $z \in \R^{256}$, and predicts the next map as a residual $\hat{s}_{t+1} = s_t + g_\theta(c_t, a_t, z)$; rollouts feed $\hat{s}_{t+1}$ back into the context, at one network evaluation per step.

\textbf{Generator.} $g_\theta$ is a ConvNeXt stack \citep{liu2022convnext} on the token grid. The context maps are concatenated along channels; the action is injected by feature-wise linear modulation (FiLM); the noise is injected by adaptive layer normalization through a four-layer mapping network, as in ROMS-IMLE \citep{vashist2026roms}; and the output head is a per-token MLP of width 1024, the ``wide head'' that FlowWM found necessary in feature space, which we attach to every method so that parameter counts match within $\pm4\%$. We use two sizes: width 256 with 8 blocks ($8.0$--$8.2$M across methods) for the synthetic benchmarks and width 384 with 12 blocks ($20.6$M) for Push-T and BAIR. The noise enters only the last three blocks (four for the larger model), so the trunk is evaluated once per transition and only the tail once per candidate; this late fusion halves training time at equal coverage (Appendix~\ref{app:impl}).

\textbf{Objective.}
\label{sec:objective}
For each transition we draw $m$ noise vectors (the \emph{candidate budget}, or budget for short), generate $m$ candidates without gradient tracking, select the one nearest to the realized outcome in token space,
\begin{equation}
j^\star = \argmin_{j \in [m]} \big\| s_{t+1} - \hat{s}^{(j)} \big\|^2 ,
\label{eq:match}
\end{equation}
and regenerate the selected candidate with gradient tracking, training it with a Geman--McClure loss \citep{geman1987statistical} on its root-mean-square distance $d$: $\rho(d) = 2u^2/(u^2+4)$, where $u = d/(2c)$ and $c$ is the running median of the matching distance. The matching is recomputed for every batch with a per-transition candidate pool \citep[as in WIMLE;][]{aghabozorgi2026wimle}, in bfloat16, without gradient tracking, and in a separate mixed-precision region (Appendix~\ref{app:impl} describes a pitfall of sharing one). The robust loss reproduced the gain reported for ROMS-IMLE in our pilot study but matters little in the final recipe (Appendix~\ref{app:impl}); the budget $m$ is the variable this paper studies. We call the resulting model \imlewm{}. Every 2{,}000 steps, a \emph{memorization monitor} draws 16 samples for each of 200 training contexts and 200 held-out contexts whose next transition is stochastic, and records the share of samples decoded onto the mode that was realized in the data. This share is $1/K$ in expectation on held-out contexts, and it exceeds $1/K$ on training contexts only if the model has learned which mode each training context took.

\textbf{Per-step matching on self-rollouts.}
\label{sec:seqmatch}
Matching whole $H$-step candidate trajectories fails: the number of trajectory modes grows like $K^H$, so the selected candidate is the one that is close to all of them on average (Appendix~\ref{app:negative}). We instead match one step at a time along the model's own rollout. Starting from a real context, at each step $h = 1,\dots,H$ we draw $m$ fresh noise vectors, select the candidate nearest to the true frame $s_{t+h}$ \emph{given the self-predicted context}, train that step, and feed the selected prediction back (detached) as the next context; with $H{=}4$ this costs $mH$ forward passes without gradient tracking and $H$ with it. Section~\ref{sec:seq} compares two alternatives for the fed-back state, a random candidate and the true frame. Selecting the candidate nearest the ground truth is the device of closest-among-top-$K$ rollouts \citep{zhang2025catk}; in \nom{} training this selection is the argmin of the loss itself.

\textbf{Inference, planning and baselines.}
\label{sec:baselines}
One predicted step is one forward pass, and $m_p$ futures for a state are obtained by batching $m_p$ noise vectors. Planning uses cross-entropy-method model-predictive control \citep[CEM-MPC;][]{rubinstein1999cem,zhou2025dinowm}, scoring each candidate action sequence by the mean, over $m_p$ imagined futures, of a distance to the goal; a variance penalty lowered the success rate even with the true simulator as the model (Appendix~\ref{app:negative}). All baselines share the backbone, wide head, context, residual parameterization and training budget unless stated otherwise: \emph{regression} (squared error); \emph{FM-50}, flow matching \citep{lipman2023flow} in the FlowWM recipe (time shift $\alpha = 13$, 50 Euler steps), and \emph{FM-2}, the same model sampled with 2 steps; \emph{iMF} \citep{geng2025meanflow,geng2025imf}, ported from the official code and trained for $3\times$ as many steps; \emph{Shortcut} \citep{frans2025shortcut,hafner2025dreamer4}, trained to the wall-clock budget of \imlewm{} with $m{=}16$; and the \emph{energy score} \citep{gneiting2007strictly} as used by \citet{makabe2026esp}, i.e.\ the \imlewm{} generator trained with the unbiased $m$-sample estimator of $\E\|Y - x\|^\beta - \tfrac{1}{2}\E\|Y - Y'\|^\beta$ ($\beta = 1$, whole-map distance, $m \in \{8, 16\}$), plus a per-token variant (Appendix~\ref{app:impl}).
